%% kapa: exam=1 round=27 year=2016 period=1 form=A subject=civil_law subject_ko=민법 seq=2 number=2
%% source: Q-Net 감정평가사 기출문제 (cst00309) artlSeq=5169402 file=extracted/1008/5169402/제27회 감정평가사 1차 문제/제27회 감정평가사 1교시문제(A형).pdf
%% conversion: kordoc 4.24.51 -> tools/gen1.py (auto); review=proofread-sonnet; flags=none
\begin{question}{2}
실종선고에 관한 설명으로 옳지 않은 것은? (다툼이 있으면 판례에 따름)
\begin{choices}
  \item 실종선고 취소의 청구를 받은 가정법원은 공시최고의 절차를 거칠 필요가 없다.
  \item 실종선고가 확정되면 실종선고를 받은 자는 실종기간이 만료한 때에 사망한 것으로 본다.
  \item 실종선고가 취소되더라도 실종기간 만료 후 실종선고 취소 전에 선의로 한 행위의 효력에는 영향을 미치지 아니한다.
  \item 실종선고의 취소가 있는 경우, 실종선고를 직접원인으로 하여 재산을 취득한 자는 선의이면 그 받은 이익이 현존하는 한도에서 반환할 의무가 있다.
  \item 부재자의 1순위 상속인이 있는 경우에 4순위의 상속인에 불과한 자는 특별한 사정이 없는 한 부재자에 대한 실종선고를 청구할 이해관계인이 될 수 없다.
\end{choices}
\end{question}
\answer{3}
